\chapter{InternVL3-2B Prompts and Output Processing}
\label{ch:appendix_d}

This appendix reproduces the prompts and output processing rules used with InternVL3-2B.
Section~\ref{sec:appendix_d_inference} reports the inference settings, Sections~\ref{sec:appendix_d_attribute_prompt} to~\ref{sec:appendix_d_binary_prompt} provide the three prompts verbatim, and Section~\ref{sec:appendix_d_parsing_rules} defines response parsing and building-level aggregation.
The prompts are preserved verbatim because changes in their wording would define a different experimental condition.

\section{Inference Settings}
\label{sec:appendix_d_inference}

Prompt wording and response constraints were developed using a fixed subset of 200 development buildings stratified by reference building type, with minority building types oversampled.
The prompts were finalised before holdout evaluation.
The model is loaded in bfloat16, and each image is processed through dynamic tiling with no more than twelve \(448 \times 448\) tiles.
Generation uses greedy decoding with sampling disabled (\texttt{do\_sample = False}) and \texttt{max\_new\_tokens = 512}.
The three highest ranked images available for each building are processed independently.

\section{Attribute Prediction Prompt}
\label{sec:appendix_d_attribute_prompt}

\begin{lstlisting}[keywords={},mathescape=false,basicstyle=\ttfamily\footnotesize,breaklines=true]
Look at this Dutch residential building street-view photo. Estimate three physical attributes of the main building in view.

Respond with EXACTLY one JSON object, no markdown fences, no extra text. You MUST provide values for ALL fields. NEVER return null.

Required keys:

- "building_type": one of
    "SFH" = single-family house, FREE-STANDING: open space visible on at
            least one side. A building sharing BOTH side walls in a
            continuous row is NEVER SFH.
    "TH"  = terraced / row house: a dwelling in a continuous row of similar
            houses, typically 2-3 storeys, in a residential street
    "AB"  = apartment block: multiple dwellings stacked in one building.
            Cues: a shared entrance with many doorbells/mailboxes, external
            access galleries, a wide flat repetitive facade, OR a narrow
            historic building of 3+ storeys in a dense city-centre street
            (these are usually subdivided into stacked apartments)
    "MFH" = collective housing without self-contained dwellings (student
            dormitory, rooming house, elderly care home)
  If unsure between TH and AB: a 2-3 storey row house in a quiet residential
  street is TH; 4+ storeys, shops at ground level, or a dense city-centre
  street means AB.

- "construction_year": integer 1800-2025. Estimate from visible evidence:
  facade style, brickwork, window shapes, roof form, detailing.

- "num_floors": integer 1-30, visible storeys above ground (count window rows).

Your JSON object must contain ALL three keys, in this order. If uncertain,
give your best single estimate - never omit a key.

Example: {"building_type": "TH", "construction_year": 1932, "num_floors": 3}
\end{lstlisting}

\section{Direct Seven-Class Prediction Prompt}
\label{sec:appendix_d_sevenclass_prompt}

The seven-class and binary prompts use the same description of visual evidence and differ only in the permitted answer space.
This keeps the visual assessment instructions consistent across the two target definitions.

\begin{lstlisting}[keywords={},mathescape=false,basicstyle=\ttfamily\footnotesize,breaklines=true]
Look at this Dutch residential building street-view photo. Estimate the building's official Dutch energy label (energielabel, NTA 8800), which rates a home's primary fossil energy use from A (most energy-efficient) to G (least energy-efficient).

Judge ONLY from visible evidence in the photo:
- window glazing: single-pane (older, less efficient) vs double / HR++ glazing
- solar panels on the roof
- signs of recent renovation or added wall insulation (new render, cladding, modern detailing)
- general age and upkeep of the facade and roof
Buildings that look newly built or recently renovated, with modern glazing or solar panels, tend toward A-B. Old, un-renovated buildings with single glazing and worn facades tend toward F-G. Most ordinary housing falls in the middle.

Respond with EXACTLY one JSON object, no markdown fences, no extra text. You MUST provide a value. NEVER return null.

Required key:
- "energy_label": exactly one of "A", "B", "C", "D", "E", "F", "G"

Example: {"energy_label": "C"}
\end{lstlisting}

\clearpage
\section{Direct Binary Prediction Prompt}
\label{sec:appendix_d_binary_prompt}

\begin{lstlisting}[keywords={},mathescape=false,basicstyle=\ttfamily\footnotesize,breaklines=true]
Look at this Dutch residential building street-view photo. Dutch homes carry an official energy label (energielabel, NTA 8800) rating primary fossil energy use from A (most energy-efficient) to G (least energy-efficient).

Decide which of two bands this building falls into:
- "A-C": the more energy-efficient band
- "D-G": the less energy-efficient band

Judge ONLY from visible evidence in the photo:
- window glazing: single-pane (older, less efficient) vs double / HR++ glazing
- solar panels on the roof
- signs of recent renovation or added wall insulation (new render, cladding, modern detailing)
- general age and upkeep of the facade and roof
Buildings that look newly built or recently renovated, with modern glazing or solar panels, tend toward A-C. Old, un-renovated buildings with single glazing and worn facades tend toward D-G.

Respond with EXACTLY one JSON object, no markdown fences, no extra text. You MUST provide a value. NEVER return null.

Required key:
- "energy_band": exactly one of "A-C", "D-G"

Example: {"energy_band": "A-C"}
\end{lstlisting}

\section{Response Parsing and Building-Level Aggregation}
\label{sec:appendix_d_parsing_rules}

Each image response is parsed before predictions are aggregated by building.
A response that fails an applicable validity rule is excluded from aggregation.
If all responses for a building are invalid, no building prediction is produced.

\begin{enumerate}
	\item The first JavaScript Object Notation object is extracted from the response. Markdown code fences and additional text before or after the object are tolerated. Common trailing comma errors are repaired before parsing.
	\item For attribute prediction, \texttt{building\_type} must be one of \{SFH, TH, MFH, AB\}, \texttt{construction\_year} must be an integer from 1800 to 2025, and \texttt{num\_floors} must be an integer from 1 to 30. All three values must be valid for the response to enter building-level aggregation.
	\item For direct energy class prediction, \texttt{energy\_label} must be one of \{A, B, C, D, E, F, G\}, and \texttt{energy\_band} must be one of \{A-C, D-G\}.
	\item Attribute predictions are aggregated by the modal building type, the median construction year, and the rounded median floor count. The implementation applies no additional rule when building type categories have equal modal counts.
	\item Direct energy class predictions are aggregated by majority vote. A tie is resolved in favour of the more energy efficient class or class group.
\end{enumerate}

These rules convert valid image responses into at most one attribute record and one direct energy class prediction for each building and target definition.
